\section{Conjugaison}

\subsection{辞書形}

La forme neutre ou forme du dictionnaire (\ltjruby{辞書形}{じしょけい}) est la forme à partir de laquelle nous répartirons puis conjuguerons les verbes.

\begin{enumerate}
    \item \textbf{五段動詞 :} Les verbes godan (\ltjruby{五段動詞}{ごだんどうし}) se terminent par un des neuf kana en う suivants :
          \begin{itemize}
              \item \ltjruby{買}{か}う
              \item \ltjruby{聞}{き}く
              \item \ltjruby{泳}{およ}ぐ
              \item \ltjruby{探}{さが}す
              \item \ltjruby{待}{ま}つ
              \item \ltjruby{死}{し}ぬ
              \item \ltjruby{飛}{と}ぶ
              \item \ltjruby{読}{よ}む
              \item なる
          \end{itemize}
          On peut noter que 死ぬ est le seul verbe du japonais moderne à se terminer par la syllabe ぬ.
    \item \textbf{一段動詞 :} Les verbes ichidan (\ltjruby{一段動詞}{いちだんどうし}) se terminent par [-i]る ou [-e]る. On ne peut pas a priori déterminer si un verbe à la forme en [-i]る ou [-e]る est du groupe 五段 ou 一段.
          \begin{itemize}
              \item \ltjruby{着}{き}る
              \item \ltjruby{食}{た}べる
          \end{itemize}
    \item \textbf{不規則動詞 :} Les verbes irréguliers (\ltjruby{不規則動詞}{ふきそくどうし}) sont des verbes très utilisés :
          \begin{itemize}
              \item する
              \item ある
              \item だ
              \item \ltjruby{行}{い}く
              \item \ltjruby{来}{く}る
          \end{itemize}
          Le verbe だ provient de la contraction de である (particule で accolée au verbe ある), et on pourra donc souvent se ramener à la conjugaison de ある, pour ensuite recontracter le résultat.
\end{enumerate}

Cette forme s'emploie à l'affirmatif neutre non-accompli :

\begin{tcolorbox}
    -- \ltjruby{毎朝}{まいあさ}、\ltjruby{六時}{ろくじ}に\ltjruby{起}{お}きる。 \ltjruby{学校}{がっこう}へ\ltjruby{行}{い}く\ltjruby{前}{まえ}に、\ltjruby{朝}{あさ}ごはんを\ltjruby{食}{た}べる。

    -- Tous les matins, je me lève à 6h. Avant d'aller à l'école, je mange un petit-déjeuner.
\end{tcolorbox}

Il existe aussi des verbes \textbf{sino-japonais}, formés d'un nom auquel on ajoute l'auxiliaire する, qui constituent un nouveau groupe qui dépend uniquement de la conjugaison du verbe irrégulier する. L'utilisation de ces verbes plutôt que leur alternative purement japonaise est considérée souvent comme plus littéraire. On peut par exemple mentionner les verbes étudier (\ltjruby{勉強}{べんきょう}する) ou se promener (\ltjruby{散歩}{さんぽ}する).

\begin{tcolorbox}
    -- \ltjruby{週末}{しゅうまつ}に\ltjruby{友達}{ともだち}と\ltjruby{買い物}{かいもの}する。

    -- Je fais du shopping avec mes amis le week-end.
\end{tcolorbox}

\subsection{未然形}

La forme inaccomplie (\ltjruby{未然形}{みぜんけい}) est communément surnommé forme en \~{}あ. Voyons comment créer cette forme pour les différents groupes :

\begin{enumerate}
    \item \textbf{五段動詞 :} On remplace le kana en [-u] par le kana en [-a] :
          \begin{itemize}
              \item 買う　>　買わ\~{}
              \item 聞く　>　聞か\~{}
              \item 泳ぐ　>　泳が\~{}
              \item 探す　>　探さ\~{}
              \item 待つ　>　待た\~{}
              \item 死ぬ　>　死な\~{}
              \item 飛ぶ　>　飛ば\~{}
              \item 読む　>　読ま\~{}
              \item なる　>　なら\~{}
          \end{itemize}
          La seule exception de cette liste sont les verbes en う, dont le kana est remplacé par わ pour une meilleur compréhension à l’oral.
    \item \textbf{一段動詞 :} On enlève la terminaison る finale :
          \begin{itemize}
              \item 着る　>　着\~{}
              \item 食べる　>　食べ\~{}
          \end{itemize}
    \item \textbf{不規則動詞 :}
          \begin{itemize}
              \item する n'a pas de forme connective particulière.
              \item ある n'a pas de forme connective particulière.
              \item だ n'a pas de forme connective particulière.
              \item 行く　>　行か\~{}
              \item \ltjruby{来}{き}る　>　\ltjruby{来}{\textcolor{couleurmatiere}{こ}}\~{}
          \end{itemize}
\end{enumerate}

\subsubsection{La négation neutre 未然形\~{}ない}

La négation neutre peut être formée aux temps accomplis et non-accomplis, par les formules suivantes :

\begin{tblr}[theme=general]{colspec={Q[1]Q[1]}}
    Accompli                                & Non accompli                          \\
    未然形\~{}\textcolor{couleurmatiere}{なかった} & 未然形\~{}\textcolor{couleurmatiere}{ない} \\
\end{tblr}

Pour les trois verbes n'ayant pas de forme connective particulière, il faudra retenir l'expression de la négation neutre.
\begin{tblr}[theme=general]{colspec={Q[1]Q[2]Q[2]}}
       & Acc.                & N.Acc.            \\
    する & しなかった               & しない               \\
    ある & なかった                & ない                \\
    だ  & で\ltjruby{は}{わ}なかった & で\ltjruby{は}{わ}ない \\
\end{tblr}



\subsection{連用形}

La forme conjonctive (\ltjruby{連用形}{れんようけい}), communément appelée forme en \~{}い, est utile notamment pour exprimer la formule polie en japonais.

\begin{enumerate}
    \item \textbf{五段動詞  :} On remplace le kana en [-u] par le kana en [-i] :
        \begin{itemize}
            \item \ruby{思}{おも}う　>　\ruby{思}{おも}い
            \item \ruby{続}{つづ}く　>　\ruby{続}{つづ}き
            \item \ruby{急}{いそ}ぐ　>　\ruby{急}{いそ}ぎ
            \item \ruby{話}{はな}す　>　\ruby{話}{はな}し
            \item \ruby{勝}{か}つ　>　\ruby{勝}{か}ち
            \item \ruby{死}{し}ぬ　>　\ruby{死}{し}に
            \item \ruby{遊}{あそ}ぶ　>　\ruby{遊}{あそ}び
            \item \ruby{休}{やす}む　>　\ruby{休}{やす}み
            \item \ruby{終}{お}わる　>　\ruby{終}{お}わり
        \end{itemize}

    \item \textbf{一段動詞 :} On enlève la terminaison る finale :
          \begin{itemize}
              \item \ruby{感}{かん}じる　>　\ruby{感}{かん}じ
              \item \ruby{答}{こた}える　>　\ruby{答}{こた}え
          \end{itemize}
    \item \textbf{不規則動詞 :}
          \begin{itemize}
              \item する　>　し
              \item ある　>　あり
              \item だ　>　であり
              \item 行く　>　行き
              \item 来る　>　\ltjruby{来}{\textcolor{couleurmatiere}{き}}
          \end{itemize}
\end{enumerate}

La forme conjonctive peut exister \textit{seule}, car elle permet de créer le nom associé à l'action d'un verbe. Par exemple avec les formes supmentionnées, 思い (la pensée), 続き (la suite), 急ぎ (la hâte), 話し (la discussion), 勝ち (la victoire), 遊び (le jeu), 休み (le repos, les vacances), 終わり (la fin), 感じ (la sensation), 答え (la réponse).

\subsubsection{La politesse 連用形\~{}ます}

Pour construire la forme polie d'un verbe, très utilisée lorsque l'on parle à un inconnu ou à une personne respectée, on utilise la forme conjonctive.

\begin{tblr}[theme = general]{colspec={Q[c]Q[l]Q[l]}}
           & Aff.       & Nég.          \\
    Acc.   & 連用形\~{}ました & 連用形\~{}ませんでした \\
    N.Acc. & 連用形\~{}ます  & 連用形\~{}ません    \\
\end{tblr}

Le verbe だ, le plus important de la langue japonaise, a un comportement particulier qu'il faut connaître :
\begin{tblr}[theme = general]{colspec={Q[c]Q[l]Q[l]}}
          & Aff. & Nég.                    \\
    Acc.  & でした  & で\ltjruby{は}{わ}ありませんでした \\
    N.Acc & です   & で\ltjruby{は}{わ}ありません    \\
\end{tblr}

\subsubsection{Le désir 連用形\~{}たい}

L'expression du désir (« je veux... ») en japonais se construit à partir de la forme nominalisée, à laquelle on ajoute le suffixe\~{}たい.

\begin{tcolorbox}
    -- \ruby{今晩}{こんばん}、私は友達と日本の映画を見たいです。

    -- Ce soir, je veux regarder un film japonais avec mon ami.
\end{tcolorbox}

Cette transformation fait en réalité du verbe un い-\textit{adjectif}, et il faut donc se ramener à leur conjugaison pour les différentes formes :

\begin{tblr}[theme = general]{colspec = {Q[c] Q[c] Q[l] Q[l]}}
         &        & Acc.             & N.Acc.        \\
    Aff. & Ntr. & 連用形\~{}たかった       & 連用形\~{}たい      \\
         & Pol.  & 連用形\~{}たかったです     & 連用形\~{}たいです    \\
         & Pol+ & 連用形\~{}たくありました    & 連用形\~{}たくあります  \\
    Nég. & Ntr. & 連用形\~{}たくなかった     & 連用形\~{}たくない    \\
         & Pol.  & 連用形\~{}たくなかったです   & 連用形\~{}たくありません \\
         & Pol+ & 連用形\~{}たくありませんでした & 連用形\~{}たくありません \\
\end{tblr}

\begin{tcolorbox}
    -- この\ltjruby{映画}{えいが}を\ltjruby{本当}{ほんとう}に\ltjruby{見}{み}たかったんですが、\ltjruby{時間}{じかん}がありませんでした。日本の\ruby{料理}{りょうり}を\ruby{食}{た}べたいです。

    -- Je voulais vraiment voir ce film, mais je n'ai pas eu le temps. Je veux manger de la cuisine japonaise.
\end{tcolorbox}

Il est très important de remarquer qu'en japonais la forme en \~{}たい est très personnelle, et donc est strictement limitée à l'expression de ses propres désirs.

\subsubsection{L'impression 連用形\~{}そう}

Pour exprimer l'idée d'impression (« Il me semble qu'il va pleuvoir » par exemple), on utilise la forme en \~{}そう ajoutée au verbe à la forme conjonctive. On forme alors un mot qui s'utilise de la même façon qu'un な-adjectif.

\begin{tcolorbox}
    -- 今にも雨が\ruby{降}{ふ}りそうです。 さっき、電車が\ruby{来}{き}そうでした。

    -- On dirait qu'il va pleuvoir à tout moment. Tout à l’heure, on aurait dit que le train allait arriver.
\end{tcolorbox}


\subsubsection{Le gérondif 連用形\~{}ながら}

Pour former le gérondif, on utilise la forme nominalisée à laquelle on ajoute\~{}ながら.

\begin{tcolorbox}
    -- \ltjruby{音楽}{おんがく}を\ltjruby{聞}{き}きながら\ltjruby{勉強}{べんきょう}するのが\ltjruby{好}{す}きです。スマホを\ltjruby{見}{み}ながら\ltjruby{運転}{うんてん}しないでください。

    -- J'aime étudier tout en écoutant de la musique. Veuillez ne pas conduire en regardant un smartphone.
\end{tcolorbox}

\subsubsection{Les tendances naturelles 連用形\~{}やすい}

Lorsque l'on ajoute やすい à la forme en い d'un verbe, on créé un adjectif variable pour décrire une action qui est facile à accomplir ou bien qu'un événement à une tendance naturelle à se produire. Cet adjectif peut donc ensuite être conjugué normalement.

\begin{tcolorbox}
    -- その本はあまり読みやすくなかったです。白い\ruby{服}{ふく}は\ruby{汚}{よご}れやすいです。この\ruby{辞書}{じしょ}は\ruby{使}{つか}いやすくて、\ruby{便利}{べんり}です。

    -- Ce livre n'était pas très facile à lire. Les vêtements blancs se salissent facilement. Ce dictionnaire est facile à utiliser et pratique.
\end{tcolorbox}

\subsubsection{Le volitif poli 連用形\~{}ましょう}

L'expression de la volonté est différente de celle du désir : elle implique une action prochaine. Elle montre que l'on cherche à entreprendre une action. C'est le fameux rôle du « let's \dots » en anglais. Les formes négative et accomplie sont peu utilisées, donc concentrons-nous sur la partie affirmative non accomplie.

Pour la forme polie (la forme neutre est créée sur une base différente), on utilise la formule
\begin{center}
    「連用形\~{}\textcolor{couleurmatiere}{ましょう}」
\end{center}

\begin{tcolorbox}
    -- 今日は\ruby{暑}{あつ}いですから、\ruby{家}{うち}で勉強しましょう。

    -- Comme il fait chaud aujourd'hui, étudions à la maison.
\end{tcolorbox}

\subsection{て形}

À l'origine, la forme en て (\ltjruby{て形}{てけい}) est créée à partir de la forme nominalisée à laquelle on ajoute le kana て. Cependant, les verbes du premier groupe ont subi des transformations avec le temps.

\begin{enumerate}
    \item \textbf{五段動詞 :} La forme en て (ou で) dépend légèrement des terminaisons.
        \begin{itemize}
            \item \ruby{言}{い}う　>　\ruby{言}{い}って
            \item \ruby{書}{か}く　>　\ruby{書}{か}いて
            \item \ruby{脱}{ぬ}ぐ　>　\ruby{脱}{ぬ}いで
            \item \ruby{殺}{ころ}す　>　\ruby{殺}{ころ}して
            \item \ruby{立}{た}つ　>　\ruby{立}{た}って
            \item \ruby{死}{し}ぬ　>　\ruby{死}{し}んで
            \item \ruby{選}{えら}ぶ　>　\ruby{選}{えら}んで
            \item \ruby{飲}{の}む　>　\ruby{飲}{の}んで
            \item \ruby{売}{う}る　>　\ruby{売}{う}って
        \end{itemize}

    \item \textbf{一段動詞 :} On enlève la terminaison る finale, et on ajoute la syllabe て.
          \begin{itemize}
              \item \ruby{見}{み}る　>　\ruby{見}{み}て
              \item \ruby{出}{で}る　>　\ruby{出}{で}て
          \end{itemize}
    \item \textbf{不規則動詞 :}
          \begin{itemize}
              \item する　>　して
              \item ある　>　あって
              \item だ　>　で
              \item 行く　>　行って
              \item \ltjruby{来}{く}る　>　\ltjruby{来}{き}て
          \end{itemize}
\end{enumerate}

\subsubsection{Énumération de verbes}

On peut, dans une même phrase, utiliser plusieurs verbes. Comme l'indique son nom, la forme suspensive ne conclut pas la phrase, et laisse la place à un autre verbe.
\begin{center}
    「\dots　て形、\dots　て形、\dots verbe」
\end{center}
\begin{tcolorbox}
    -- 彼は毎晩お茶を\ruby{飲}{の}んで、漢字を\ruby{習}{なら}って、シャワーを\ruby{浴}{あ}びる。

    -- Tous les soirs, il boit du thé, apprend des kanjis et prend une douche.
\end{tcolorbox}

Pour un même thème, l'ordre d'appparition des verbes doit se faire dans l'ordre chronologique, et une proposition suivante peut être impliquée par la précédente. Toutefois, en précisant le thème à chaque fois, on peut parler d'actions simultanées réalisées par des sujets différents. Le caractère accompli ou non de chaque verbe est signifié par le dernier verbe.

Il existe, de même qu'avec les noms avec la particule や, un moyen d'énumérer une liste non exhaustive de verbes, par la formule
\begin{center}
    「\dots　たり形、\dots　verbe-たり形、\dots verbe-たり形\~{}する」
\end{center}
Cette liste n'implique pas d'ordre chronologique particulier, et le caractère accompli se marque par la conjugaison de する.
\begin{tcolorbox}
    -- 私は\ruby{妹}{いもうと}とケーキを{作}{つく}ったり、庭で自転車を\ruby{直}{なお}したり、\ruby{髪}{かみ}を\ruby{石鹸}{せっけん}で\ruby{洗}{あら}ったりしながら、\ruby{一日中}{いちにちじゅ}ラジオを聞いて。

    -- J'ai écouté la radio toute la journée, en faisant un gâteau avec ma petite sœur, en réparant mon vélo dans le jardin, ou encore, en lavant mes cheveux au savon.
\end{tcolorbox}

Si l'on veut que l'un des verbes énumérés soient au négatif, on peut utiliser la forme suspensive d'un verbe négatif : on remplace le ない du négatif par なくて.

\subsubsection{Le progressif て形\~{}いる}

Pour exprimer la progressivité (« en train de \dots »), on utilise la forme en て associée au verbe いる :

\begin{center}
    「sujet　が　て形\~{}いる」
\end{center}

\begin{tcolorbox}
    -- 今、本を読んでいます。

    -- Je lis un livre en ce moment.
\end{tcolorbox}

On emploie aussi cette forme pour décrire un état qui demeure après qu'un changement a eu lieu (avec un verbe décrivant un changement soudain).

\begin{tcolorbox}
    -- 田中さんは\ruby{結婚}{けっこん}していて、東京に住んでいます

    -- M. Tanaka est marié et habite à Tokyo.
\end{tcolorbox}

Cette forme s'emploie aussi lorsque l'on s'exprime sur \textit{des habitudes ou des occupations}, et qu'une notion de régularité sur le long terme intervient.

\begin{tcolorbox}
    -- 毎日日本語お勉強しています。

    -- J'étudie le japonais tous les jours.
\end{tcolorbox}

\subsubsection{La demande て形\~{}ください}

La forme てください est la forme qui permet d'exprimer une demande ou un ordre de façon polie.
\begin{center}
    「\dots　て形　+　\~{}ください」
\end{center}
aura pour signification « pouvez-vous \dots, s'il vous plaît ».
\begin{tcolorbox}
    -- \ruby{暑}{あつ}いですから、\ruby{窓}{まど}を\ruby{開}{あ}けてください。

    -- S'il te plaît, ouvre la fenêtre parce qu'il fait chaud.
\end{tcolorbox}

\subsubsection{La permission て形\~{}もいい}

Cette forme, qui signifie littéralement « même si tu fais verbe, c'est bien », permet de donner ou de demander une permission. On peut traduire effectivemet traduire « même si tu manges, c'est bien » par « tu peux manger », bien plus naturel en français.

\begin{tcolorbox}
    -- ピアノを\ruby{弾}{ひ}いてもいいですか。はい、どうぞ。

    -- Puis-je jouer du piano ? Oui, je vous en prie.
\end{tcolorbox}

\subsubsection{L'interdiction て形\~{}はいけません}

Pour formuler une interdiction forte, comme dans une loi ou des régulations, on emploie la forme て形\~{}はいけません.

\begin{tcolorbox}
    -- ここで写真を\ruby{撮}{と}ってはいけません。

    -- Vous ne devez pas prendre de photo ici.
\end{tcolorbox}

\subsection{仮定形}

La forme hypothétique (\ltjruby{仮定形}{かていけい}) se forme de la façon suivante :

\begin{enumerate}
    \item \textbf{五段動詞 :} On remplace le kana en [-u] par le kana en [-e].
          \begin{itemize}
              \item \ltjruby{使}{つか}う　>　\ltjruby{使}{つか}え\~{}
              \item \ltjruby{歩}{ある}く　>　\ltjruby{歩}{ある}け\~{}
              \item \ltjruby{泳}{およ}ぐ　>　\ltjruby{泳}{およ}げ\~{}
              \item \ltjruby{消}{け}す　>　\ltjruby{消}{け}せ\~{}
              \item \ltjruby{持}{も}つ　>　\ltjruby{持}{も}て\~{}
              \item \ltjruby{死}{し}ぬ　>　\ltjruby{死}{し}ね\~{}
              \item \ltjruby{呼}{よ}ぶ　>　\ltjruby{呼}{よ}べ\~{}
              \item \ltjruby{休}{やす}む　>　\ltjruby{休}{やす}め\~{}
              \item \ltjruby{取}{と}る　>　\ltjruby{取}{と}れ\~{}
          \end{itemize}
    \item \textbf{一段動詞 :} On remblace la syllabe る par れ.
          \begin{itemize}
              \item \ltjruby{起}{お}きる　>　\ltjruby{起}{お}きれ\~{}
              \item \ltjruby{寝}{ね}る　>　\ltjruby{寝}{ね}れ\~{}
          \end{itemize}
    \item \textbf{不規則動詞 :}
          \begin{itemize}
              \item する　>　すれ\~{}
              \item ある　>　あれ\~{}
              \item だ　>　であれ\~{}
              \item 行く　>　行け\~{}
              \item 来る　>　来れ\~{}
          \end{itemize}
\end{enumerate}

\subsubsection{La forme conditionnelle 仮定形\~{}ば}

La forme conditionnelle, comme son nom l'indique, permet principalement d'établir des énoncés de condition logique.

\begin{tcolorbox}
    -- ボタンを\ruby{押}{お}せば、ドアが\ruby{開}{あ}きます。

    -- Si vous appuyez sur le bouton, la porte s'ouvrira.
\end{tcolorbox}

Elle permet aussi, dans le même style, d'énoncer des vérités générales ou des proverbes.

\begin{tcolorbox}
    -- \ruby{春}{はる}になれば、\ruby{花}{はな}が\ruby{咲}{さ}きます。

    -- Quand le printemps arrivera, les fleurs écloront.
\end{tcolorbox}

De façon un peu différente, cette forme permet aussi de demander ou donner des conseils.

\begin{tcolorbox}
    -- どうすればいいですか。

    -- Que devrais-je faire ?
\end{tcolorbox}

\subsection{た形}

Pour créer la forme en た d'un verbe (た形, たけい), on imite la création de la forme en て, en remplaçant respectivement て par た et で par だ.

\begin{enumerate}
    \item \textbf{五段動詞 :}
          \begin{itemize}
              \item 言う　>　言った
              \item 書く　>　書いた
              \item 脱ぐ　>　ぬいで
              \item 殺す　>　殺した
              \item たつ　>　立った
              \item 死ぬ　>　死んで
              \item 選ぶ　>　選んで
              \item 飲む　>　飲んで
              \item 売る　>　うった
          \end{itemize}
    \item \textbf{一段動詞 :} On enlève la terminaison る finale, et on ajoute la syllabe た.
          \begin{itemize}
              \item 見る　>　見た
              \item 出る　>　出た
          \end{itemize}
    \item \textbf{不規則動詞 :}
          \begin{itemize}
              \item する　>　した
              \item ある　>　あった
              \item だ　>　だった
              \item 行く　>　行った
              \item \ltjruby{来}{く}る　>　\ltjruby{来}{き}た
          \end{itemize}
\end{enumerate}

    En employant la forme ainsi faite, on exprime les verbes à la forme accomplie neutre.

    \begin{tcolorbox}
        -- 朝ご飯を食べた。

        -- J'ai pris mon petit-déjeuner.
    \end{tcolorbox}

    \subsection{意向形}

    La forme volitive neutre (\ruby{意向形}{いこうけい}) permet d’exprimer la volonté d’une action prochaine, qui peut se traduire par « let’s\dots » en anglais. Pour la former :

    \begin{enumerate}
        \item \textbf{五段動詞 :} On remplace le kana en [-u] par le kana en [-o], auquel on ajoute う.
        \item \textbf{一段動詞 :} On enlève la terminaison る finale, et on ajoute le suffixe\~{}よう.
        \item \textbf{不規則動詞 :}
          \begin{itemize}
              \item する　>　しよう
              \item ある　>　あろう
              \item だ n'a pas de forme volitive neutre
              \item 行く　>　行こう
              \item 来る　>　\ruby{来}{こ}よう
          \end{itemize}
    \end{enumerate}

    \begin{tcolorbox}
        -- 明日は学校が休みだから、友達と\ruby{一緒}{いっしょ}に映画を見に行こう。

        -- Demain, comme il n'y a pas école, allons voir un film avec nos amis.
    \end{tcolorbox}






